% GUID v6, Appendix: the readings that add no prediction, in full (moved from Part I).
\section{Readings in detail}\label{app:readings}

Section~\ref{sec:readings} lists the readings that add no prediction a measurement could tell apart from the established account. This appendix gives each in full.
\subsection{The thermodynamics of the frame}\label{sec:thermo}

\subsubsection{The existence of a boundary is its temperature}

GUID defines the existence a boundary asserts as the product of the local strength of gravity $g$ and Planck's constant $h$, divided by
the speed of light:
\begin{equation}\label{eq:Eex}
E_{\text{ex}}\equiv\frac{g\,h}{c}.
\end{equation}

\begin{proposition}[existence and the Unruh temperature]\label{pr:unruh}
$E_{\text{ex}}$ is the only energy that an acceleration, Planck's constant and the speed of light can form, and
\begin{equation}\label{eq:unruh}
E_{\text{ex}}=4\pi^2\,k_BT_U,\qquad k_BT_U=\frac{\hbar g}{2\pi c},
\end{equation}
where $T_U$ is the temperature an observer of proper acceleration $g$ assigns to the vacuum (Unruh~\cite{unruh}). At a horizon of surface
gravity $\kappa$ it is $4\pi^2$ times the Hawking temperature $k_BT_H=\hbar\kappa/2\pi c$; for a Schwarzschild hole $\kappa=c^4/4GM$.
\end{proposition}

\begin{proof}
With dimensions $[g]=LT^{-2}$, $[h]=ML^2T^{-1}$, $[c]=LT^{-1}$, the monomial $g^ah^bc^k$ has the dimension of energy, $ML^2T^{-2}$, only
for $(a,b,k)=(1,1,-1)$ \chk{Q1}. With $h=2\pi\hbar$, $E_{\text{ex}}/k_BT_U=4\pi^2$.
\end{proof}

The uniqueness means the form of~\eqref{eq:Eex} is not itself evidence; quantum field theory supplies the factor and the meaning. At standard
gravity $T_U=4.0\times10^{-20}$\,K; at the horizon of a solar-mass hole $T_H=6.2\times10^{-8}$\,K \chk{Q1}. \emph{Verdict: holds as a
reading} of the Unruh and Hawking temperatures.

\subsubsection{Time: the lapse}

In the GUID reading a clock asserts its existence at a rate that falls where the existence field is strong. In general relativity that
rate is the lapse, $N=\sqrt{1+2\Phi/c^2}$ in a weak field and $N=\sqrt{1-2GM/rc^2}$ outside a mass, which vanishes at the horizon
\chk{Q2}. It reproduces the redshift over the 22.5\,m Pound--Rebka tower, $2.46\times10^{-15}$~\cite{pr60}, and the rate difference of two
optical clocks 33\,cm apart, $3.6\times10^{-17}$~\cite{chou}. \emph{Verdict: holds}; it adds nothing to general relativity.

\subsubsection{Einstein's equation as an equation of state}

Jacobson derived Einstein's equation from thermodynamics~\cite{jacobson95}, and GUID's two thermal inputs are exactly his. Take any point
$p$ and any local Rindler horizon through it, generated by a null vector field $k^a$ with affine parameter $\lambda$, vanishing expansion and
shear at $p$, and approximate boost Killing field $\chi^a=-\kappa\lambda k^a$. Assume:
\begin{enumerate}[nosep]
\item the horizon has the Unruh temperature $T=\hbar\kappa/2\pi$ (units $c=k_B=1$), Proposition~\ref{pr:unruh};
\item its entropy is proportional to its area, $S=sA$, the boundary's floor (Hypothesis~\ref{h:horizon});
\item the Clausius relation $\delta Q=T\,\dd S$ holds for the energy flux across it.
\end{enumerate}
The heat flux is $\delta Q=\int T_{ab}\chi^a\dd\Sigma^b=-\kappa\int\lambda\,T_{ab}k^ak^b\,\dd\lambda\,\dd\mathcal A$. The area change follows
from the Raychaudhuri equation, $\dd\theta/\dd\lambda=-R_{ab}k^ak^b$ to first order, as $\delta A=-\int\lambda R_{ab}k^ak^b\,\dd\lambda\,\dd\mathcal A$.
Equating $\delta Q=(\hbar\kappa/2\pi)s\,\delta A$ for every $p$ and every null $k^a$ gives
$R_{ab}+f g_{ab}=(2\pi/\hbar s)T_{ab}$ for some function $f$; the contracted Bianchi identity and $\nabla^aT_{ab}=0$ fix
$f=-R/2+\Lambda$, so
\begin{equation}\label{eq:jacobson}
G_{ab}+\Lambda g_{ab}=\frac{2\pi}{\hbar s}\,T_{ab},\qquad G=\frac{1}{4\hbar s}\ \ (c=1).
\end{equation}
Newton's constant is set by the entropy per area, $s=1/4L_P^2$ \chk{Q1}, and the cosmological constant enters as an integration
constant. Jacobson's own conclusion fits GUID's classical frame: on this view, ``it may be no more appropriate to quantize the Einstein
equation than it would be to quantize the wave equation for sound in air''~\cite{jacobson95}. Two limits: DMP supplies neither the value of $\eta$ nor
a manifold, so the derivation is borrowed, not produced; and the entanglement-equilibrium version~\cite{jacobson16} needs a state, not a
schedule, which the axioms do not define. \emph{Verdict: holds as a reading}; the constant is open.

\subsubsection{Information radiation and the floor}

DMP's T31 says that a containing boundary emits only thermal radiation, carrying ``entropy and no information'', and that an outside reader
learns of the inside only its size and its staleness. With Hypothesis~\ref{h:horizon} that is the semiclassical horizon: a black hole of mass
$M$ radiates at $T_H$ and an outside observer knows its mass, spin and charge \tagP. Erasing a bit costs at least $k_BT\ln2$ of heat
(Landauer~\cite{landauer}); in the GUID reading entropy is the inefficiency of converting between existence and information, which is
Landauer's principle in other words. \emph{Verdict: holds as a reading.}

\subsection{The containment split gives Christodoulou's mass formula}

DMP splits a contained boundary into what is inside and the counts every outside reader sees \tagD. Give the contained part a mass
$M_{\text{irr}}$ that sets a Schwarzschild boundary of radius $2GM_{\text{irr}}/c^2$; let the spin count $J$ be carried at that boundary as
a momentum $p=Jc^2/(2GM_{\text{irr}})$ (an identification, \tagH); and combine the two as energy and momentum combine in special relativity. In
units $G=c=1$,
\begin{equation}\label{eq:christo}
M^2=M_{\text{irr}}^2+\left(\frac{J}{2M_{\text{irr}}}\right)^2 .
\end{equation}
This is Christodoulou's mass formula for the Kerr black hole~\cite{christodoulou}, with $M_{\text{irr}}=\sqrt{A/16\pi}$ \tagP; solved for
$M_{\text{irr}}$ it returns Kerr's horizon area for every spin $a=J/M^2\in[0,1]$, has no real solution above $a=1$ (the spin bound), and
gives the extractable share $1-M_{\text{irr}}/M=1-1/\sqrt2\approx29\%$ at maximal spin \chk{Q10}. The derivation uses DMP's split, two
physical inputs and one identification, and gives no number general relativity does not. \emph{Verdict: holds as a rederivation.}

\subsection{The universe at its own Schwarzschild threshold}

For a flat Friedmann universe with Hubble rate $H$, the mass inside the Hubble radius, $M_H=\tfrac43\pi(c/H)^3\rho_c$ with
$\rho_c=3H^2/8\pi G$, has a Schwarzschild radius equal to the Hubble radius,
\begin{equation}\label{eq:kugel}
\frac{2GM_H}{c^2}=\frac cH ,
\end{equation}
exactly and at every epoch \chk{Q12}. This is the coincidence behind Pathria's ``universe as a black hole''~\cite{pathria}. GUID reads the
early universe as a Kugelblitz, a horizon formed by energy alone, that could not contain its own existence; the reading holds in the sense
of~\eqref{eq:kugel} and at every epoch, not only the earliest, and the universe is not a black hole in the sense of general relativity,
having no outside. \emph{Verdict: holds as a reading.}

\subsection{Inflation as the possibility space of existence}

In GUID, inflation is the space that existence makes for itself wherever existence is possible. Physically it is a period of nearly
exponential expansion, $a\propto e^{Ht}$, which dilutes a spatial curvature $|\Omega-1|$ by $e^{-2N}$ over $N$ e-folds (Guth~\cite{guth})
\tagP. Afterwards the curvature share grows again, as $a^2$ while radiation dominates and as $a$ while matter does; from a reheating
temperature of $10^{15}$\,GeV to today it grows by about $5\times10^{51}$, so an order-one initial curvature needs about 62 e-folds to meet
Planck's bound $|\Omega_k|\lesssim5\times10^{-3}$~\cite{planck18} \chk{Q13}. The newest combination, DESI's second data release with the
microwave background, finds $10^3\,\Omega_K=2.3\pm1.1$, a $2.1\sigma$ departure from flatness that inflation and $\Lambda$CDM share and that
DESI does not claim as a detection (Section~\ref{sec:data}, \chk{data:D3.7}). \emph{Verdict: holds as a reading}; it supplies the motive,
not the number.

\subsection{The VIBE coefficient}

GUID names $\alpha$ the \emph{VIBE coefficient}, after the Variable Influence Balancing Equation~\eqref{eq:vibe}, and reads it as the strength
with which existence interacts with itself before it must re-ground itself. In quantum electrodynamics $\alpha$ is the strength with which
a charge interacts with the electromagnetic field and so, through it, with itself; each emission and reabsorption of a virtual photon costs
one power of $\alpha$. The cleanest measure is the electron's anomalous magnetic moment, whose first term $a_e=\alpha/2\pi$
(Schwinger~\cite{schwinger}) gives the measured value~\cite{fan} to $0.15\%$ \chk{Q16}. The same constant is the ratio of the electron's
three lengths, $r_e/\bar\lambda_C=\bar\lambda_C/a_0=\alpha$ \chk{Q16}. \emph{Verdict: holds as a reading}; the VIBE coefficient is a name for
$\alpha$, not a new constant.

